%% SOURCE: marquez-zavalaDatabase15000Strain2026 in the library mirror (paper.md sha256 952805b4205164dbf43d350eacf48e1c02f9aaad8fe9bc5529d4da358e1d8c68, DOI 10.1016/j.ymben.2026.03.017); free-fatty-acid counts from experiments/008-xue-ffa/results/ffa_epistatic_path_accessibility.csv; Kuzmin corpus counts from experiments/010-kuzmin-tmi/scripts/kuzmin_ladder_census.py; inference_4 strata from experiments/010-kuzmin-tmi/scripts/inference_4_panel_design.py

\section{Why regulation acting on metabolism\secstatus{todo}}

The \texttt{inference\_4} axes were named before this argument was written, so the
argument is owed rather than assumed. Four claims support it, three of them
measured and one of them explicitly a hypothesis.

\subsection{Strain design converges on a small set of metabolic genes\secstatus{todo}}

Marquez-Zavala 2026 (DOI \texttt{10.1016/j.ymben.2026.03.017}) text-mined over
15{,}000 strain-design articles spanning 25 years, extracted gene-level
interventions from 8{,}885 of them, and normalized every gene to a KEGG ortholog
and a BiGG reaction. \org{S. cerevisiae} and \org{E. coli} together account for
34.6 percent of the database, and the target genes concentrate sharply:

\begin{quote}\footnotesize
``We can observe that the vast majority of frequently targeted genes are
associated with reactions in central carbon metabolism (Fig. 4). Not
surprisingly, the most frequently targeted area revolves around the pyruvate
node, a key branching point of metabolism, distributing flux towards respiratory
and fermentative pathways. \dots{} Another commonly targeted area is the
branching point between upper glycolysis and the pentose-phosphate pathway,
which can control the production of NADPH, an essential cofactor for
biosynthetic pathways, and is an entry point for the production of aromatic
amino acids, which are precursors of many commercially relevant compounds.''
\end{quote}

Their sharper finding is about direction rather than identity: ``The most
distinguishing feature among strain design strategies seems not to be which genes
are targeted, but rather the direction in which they are modified (increased or
decreased expression).'' So the field returns to the same few dozen loci across
products and across hosts, and turns them up or down.

Two consequences follow. A fitness liability created at one of those loci is not
a one-off, it is a recurring cost paid again by every campaign that touches the
same branch point. And the pathway areas the survey names are enumerable in
yeast-GEM, so ``a gene strain engineering actually deletes'' can be defined by
subsystem membership rather than by hand. That definition is what the
\term{engineering axis} means in the next section: the roster genes yeast-GEM
9.0.2 assigns to glycolysis and gluconeogenesis, the pentose phosphate pathway,
the citric acid cycle, or pyruvate metabolism.

\subsection{Rechanneling flux costs growth, and growth is the binding constraint\secstatus{todo}}

Deleting a branch to redirect carbon removes a redox sink or a precursor drain
that the cell was using. In a fed-batch process the growth rate sets volumetric
productivity, so a design that raises yield while lowering growth can lose on
titer per hour. The question a strain engineer faces is therefore not only
whether a deletion redirects flux, but what restores the growth it cost.

\subsection{The rescue is plausibly regulatory, and it cannot be additive\secstatus{todo}}

\textbf{Hypothesis (untested here):} when a flux-rechanneling deletion leaves the
regulatory network still driving the removed branch, or failing to relieve a now
limiting one, deleting the responsible regulator restores growth without
restoring the competing flux.

The structural half of that hypothesis is already measured. The free-fatty-acid
panel of \protect\file{experiments/008-xue-ffa} is an in-house unpublished design
on the standard three-deletion chassis \gene{pox1}$\Delta$ \gene{faa1}$\Delta$
\gene{faa4}$\Delta$, whose chassis anchor is DOI
\protect\file{10.1016/j.ymben.2013.07.003}, crossed combinatorially against ten
transcription-factor deletions. Its best triple reaches $2.045$ times the chassis
titer. Regulator deletions stacked on a flux chassis therefore produce large
effects, and the effects are not additive.

What makes such an effect invisible to any lower-order experiment is that neither
part shows it alone. The regulator deletion by itself does nothing, because the
regulon it controls is not yet stressed, and the chassis deletions by themselves
are costly. Only the combination pays. That is the quantity $\tau$ measures, and
it is the reason a trigenic predictor is the relevant tool rather than a pair of
digenic screens.

\subsection{The combination cannot be reached by walking to it\secstatus{todo}}

Two independent measurements say a greedy stack-and-keep search does not arrive
at these strains, and both are already in this document.

\begin{table}[H]
\centering
\footnotesize
\begin{tabular}{llrr}
\toprule
Screen & Readout & Reachable greedily & Total \\
\midrule
Kuzmin trigenic corpus & growth, beats wild type & 12{,}267 (27.6\%) & 44{,}462 \\
Free-fatty-acid panel  & titer, monotone path    & 2 (1.7\%)         & 120 \\
\bottomrule
\end{tabular}
\caption{How much of the win a greedy search can see. In the Kuzmin corpus
32{,}195 of the triples that beat wild type sit behind a step that lowers
fitness. In the free-fatty-acid panel every one of the 720 deletion orders is
measured rather than predicted, and only 2 of them never step below the chassis;
allowing each step 1 standard error of slack raises that to 9 of 120 triples.
Sources are \protect\file{experiments/010-kuzmin-tmi/scripts/kuzmin_ladder_census.py}
and \protect\file{experiments/008-xue-ffa/results/ffa_epistatic_path_accessibility.csv}.}
\end{table}

The free-fatty-acid result is the stronger of the two, because the destination is
real and large while the route is not: the best triple sits at twice the chassis
titer, and 648 of the 720 orders drop below the chassis at the first deletion.
The endpoint has to be predicted, since no sequence of local improvements finds
it.

\subsection{Why cross the two axes rather than stay inside one\secstatus{todo}}

An all-metabolic triple mostly restates flux balance, which a constraint-based
model already supplies without a learned predictor, and the previous space says
the model has nothing to add there either: not one of the 1{,}183{,}337
all-metabolic \texttt{inference\_1} triples clears the positive call, their best
prediction being $+0.073$. An all-regulatory triple is not a strain design.

The cross is also close to a genuine partition, 15 of the 934 roster genes
carrying both roles, so a triple mixes two kinds of biology rather than three
members of one pathway. And it is the shape a design already has, a metabolic
chassis with regulatory tuning on top.

\subsection{What this panel can and cannot show\secstatus{todo}}

The assay reads colony growth on rich medium. It can therefore test whether a
regulator deletion rescues the growth cost of flux-rechanneling deletions, the
recurring cost every campaign at these loci pays. It cannot test titer. The titer
half of the argument rests on the free-fatty-acid panel, a different screen in a
different background with a different readout.
